{% raw %}
\documentclass{article}
\usepackage{spconf}
\usepackage{blindtext, graphicx}
\usepackage{bm}
\usepackage{booktabs}
\usepackage{subfig}
\usepackage[font=footnotesize]{caption}
\usepackage{hyperref}


\usepackage{localmath}

\newcommand{\mnote}[1]{\marginpar{\tiny{#1}}}

%%%%%%%%%%% DANGER %%%%%%%%%%%%%%%%%
%\addtolength{\columnsep}{-0.1mm}
%%%%%%%%%%%%%%%%%%%%%%%%%%%%%%%%%%%%


% correct bad hyphenation here
\hyphenation{op-tical net-works semi-conduc-tor}

% measure text and linewidth
% \usepackage{layouts}

\graphicspath{{figures/}}

\begin{document}
\ninept

{% endraw %}

\title{{ "{" }}{{cookiecutter.title}}{{ "}" }}

% Single address.
% ---------------
\name{{"{"}}{{cookiecutter.author}}%
\thanks{{"{"}}
The research presented in this paper is reproducible. Code and data are available at \protect\url{http://github.com/{{cookiecutter.code_repo}}{{ "}" }}.%
{{ "}}" }}
\address{{"{"}}{{cookiecutter.affiliation}}{{"}"}}

{% raw %}

% For example:
% ------------
%\address{School\\
%	Department\\
%	Address}
%
% Two addresses (uncomment and modify for two-address case).
% ----------------------------------------------------------
% \twoauthors
%  {A. Robin Scheibler, B. Author-two\sthanks{Thanks to XYZ agency for funding.}}
% 	{School A-B\\
% 	Department A-B\\
% 	Address A-B}
%  {C. Author-three, D. Author-four\sthanks{The fourth author performed the work
% 	while at ...}}
% 	{School C-D\\
% 	Department C-D\\
% 	Address C-D}


\maketitle

\begin{abstract}
  Summarize the paper here.
\end{abstract}
\begin{keywords}%
About, five, key, words, here.
\end{keywords}

% To display textwidth and linewidth, uncomment package layouts 
% just above begin{document} and the two following lines
% textwidth: \printinunitsof{in}\prntlen{\textwidth}
% linewidth: \printinunitsof{in}\prntlen{\linewidth}

\section{Introduction}

Don't you forget to mention all these related works.
Reference to test \ffref{test}.

\begin{figure}
    \centering
    \includegraphics[width=.7\linewidth]{test.pdf}
    \caption{A test figure}
    \label{fig:setup}
    \vspace{-0mm}
\end{figure}

\section{Model}

Describe the background, notation, and signal model here.

\section{Algorithm}

Now comes the part about the algorithms.

\section{Experiments}

Why not show the results in a table like \tref{gamma} ?
It is likely that we use \texttt{pyroomacoustics} for something \cite{scheibler2017pyroomacoustics}.

\begin{table}
    \centering
    \begin{tabular*}{\linewidth}{@{\extracolsep{\fill}}lccccccccc@{}}
        \toprule
         &       &          & \multicolumn{7}{c}{{\footnotesize Number of echoes $K$}} \\
         \cmidrule{4-10}
         & anechoic & learn & 0 & 1 & 2 & 3 & 4 & 5 & 6 \\
         \cmidrule{2-10}
         $\gamma = $ & $10$ & $10^{-1}$ & $10$ & $10^{-3}$ & 0 & 0 & 0 & 0 & 0 \\
         \bottomrule
    \end{tabular*}
    \caption{A great table}
    \tlabel{gamma}
    \vspace{-0mm}
\end{table}


\begin{figure*}
    \centering
    \subfloat[mu_univ][Test left]{\includegraphics[width=0.45\linewidth]{test.pdf}\label{fig:mu_univ}}
    \hfill
    \subfloat[mu_spkr][Test right]{\includegraphics{test.pdf}\label{fig:mu_spkr}}
    \caption{A figure with two subfigures}
    \flabel{results}
    \vspace{-0mm}
\end{figure*}


% Uncomment the following to push the conclusion to the end of the page
%\vfill

\section{Conclusion}

Conclude the paper here

\clearpage

\bibliographystyle{IEEEtran}
\bibliography{refs}

\end{document}
{% endraw %}
